\documentclass[11pt]{article}

\usepackage[margin=1in]{geometry}
\usepackage{amsmath}
\usepackage{amssymb}
\usepackage{bm}
\usepackage{mathtools}

\newcommand{\dd}{\mathrm{d}}
\newcommand{\Div}{\nabla\!\cdot}
\newcommand{\Grad}{\nabla}
\newcommand{\Er}{E_{\mathrm r}}
\newcommand{\Fr}{\bm{F}_{\mathrm r}}
\newcommand{\FrD}{\bm{F}_{\mathrm D}}
\newcommand{\Prad}{\mathsf{P}_{\mathrm r}}
\newcommand{\Gmom}{\bm{G}}
\newcommand{\Gzero}{G^0}
\newcommand{\ShatE}{S_E}
\newcommand{\ShatF}{\bm{S}_F}
\newcommand{\chat}{\hat{c}}
\newcommand{\Erzero}{E_{\mathrm r,0}}
\newcommand{\SErsla}{S_{E,\hat{c}}}
\newcommand{\SFrsla}{\bm{S}_{F,\hat{c}}}
\newcommand{\Ddiff}{D_{\mathrm r}}

\title{Manufactured Radiation Sources for a Taylor--Green--Brunner Radiation-Hydrodynamics Test}
\author{}
\date{}

\begin{document}

\maketitle

\section{Purpose}

The MAPP document describes a Taylor--Green--Brunner radiation-hydrodynamics
test in which a Taylor--Green vortex velocity is prescribed, the radiation
source cancels the hydrodynamic pressure force, the density remains constant,
and the velocity remains unchanged.  This note derives the continuum
manufactured sources needed to realize that problem.  The MAPP text places the
test in a multirate setting where radiation diffusion is one of the slow
implicit physics packages, so the main derivation below is for a grey
radiation diffusion equation.  A P1 two-moment derivation is retained as a
comparison because it is the closest form to Quokka's radiation transport
module.

The essential point is that the gas momentum equation is affected by radiation
through the radiation four-force \((\Gzero,\Gmom)\).  An external source added
only to the radiation energy equation cannot by itself cancel a gas pressure
force.  In a diffusion code this coupling is simpler: the radiation pressure
is \(p_{\mathrm r}=\Er/3\), so the gas momentum equation contains the gradient
of the total pressure \(p+\Er/3\).

\section{Quokka radiation-hydrodynamics equations}

Ignoring gravity, magnetic fields, dust, chemistry, and optically thin cooling,
the relevant full-speed \((\chat=c)\) equations are
\begin{align}
  \partial_t \rho + \Div(\rho \bm{v}) &= 0, \label{eq:mass}\\
  \partial_t(\rho \bm{v}) + \Div(\rho \bm{v}\bm{v} + p\mathsf{I}) &= \Gmom, \label{eq:mom}\\
  \partial_t \Er + \Div \Fr &= -c\Gzero + \ShatE, \label{eq:rad-energy}\\
  \partial_t \Fr + \Div(c^2\Prad) &= -c^2\Gmom + \ShatF. \label{eq:rad-flux}
\end{align}
For a P1 closure,
\begin{equation}
  \Prad = \frac{\Er}{3}\mathsf{I},
  \qquad
  \Div(c^2\Prad) = \frac{c^2}{3}\Grad \Er.
  \label{eq:p1}
\end{equation}
Therefore, for any desired exact radiation fields
\((\Er^\star,\Fr^\star)\) and exact radiation four-force
\((G^{0,\star},\Gmom^\star)\), the required radiation moment sources are the
residuals
\begin{align}
  \ShatE^\star
    &= \partial_t \Er^\star + \Div \Fr^\star + cG^{0,\star},
    \label{eq:general-energy-source}\\
  \ShatF^\star
    &= \partial_t \Fr^\star + \frac{c^2}{3}\Grad \Er^\star
       + c^2\Gmom^\star.
    \label{eq:general-flux-source}
\end{align}
These signs follow from moving the Quokka matter-radiation exchange terms,
\(-c\Gzero\) and \(-c^2\Gmom\), to the left-hand side of the radiation
equations.

\section{Taylor--Green vortex fields}

Use a two-dimensional periodic domain, with the solution independent of \(z\).
Let
\begin{equation}
  \bm{v}_{\mathrm{TG}}(x,y)
  =
  \begin{pmatrix}
    U\sin(kx)\cos(ky)\\
   -U\cos(kx)\sin(ky)\\
    0
  \end{pmatrix},
  \qquad
  \rho^\star = \rho_0.
  \label{eq:tgv}
\end{equation}
This velocity is exactly solenoidal:
\begin{equation}
  \Div \bm{v}_{\mathrm{TG}}
  = Uk\cos(kx)\cos(ky) - Uk\cos(kx)\cos(ky) = 0.
  \label{eq:tgv-divergence}
\end{equation}

For later use, define the canonical Taylor--Green pressure perturbation
\begin{equation}
  \Pi_{\mathrm{TG}}(x,y)
  =
  \frac{\rho_0 U^2}{4}
  \left[\cos(2kx)+\cos(2ky)\right].
  \label{eq:tg-pressure}
\end{equation}
The advective acceleration is
\begin{equation}
  \rho_0(\bm{v}_{\mathrm{TG}}\cdot\Grad)\bm{v}_{\mathrm{TG}}
  =
  \frac{\rho_0 U^2 k}{2}
  \begin{pmatrix}
    \sin(2kx)\\
    \sin(2ky)\\
    0
  \end{pmatrix}
  = -\Grad \Pi_{\mathrm{TG}}.
  \label{eq:tgv-advective-acceleration}
\end{equation}

\section{Radiation force required to keep the velocity unchanged}

Equation~\eqref{eq:mom} can be written in primitive form as
\begin{equation}
  \rho\left(\partial_t\bm{v} + \bm{v}\cdot\Grad\bm{v}\right) + \Grad p
  = \Gmom.
  \label{eq:primitive-momentum}
\end{equation}
Thus the most general manufactured radiation force that makes a prescribed
velocity field \(\bm{v}^\star\) an exact solution is
\begin{equation}
  \Gmom^\star
  =
  \rho^\star
  \left(\partial_t\bm{v}^\star + \bm{v}^\star\cdot\Grad\bm{v}^\star\right)
  + \Grad p^\star.
  \label{eq:general-force}
\end{equation}
For the steady Eulerian Taylor--Green field in Eq.~\eqref{eq:tgv},
\(\partial_t\bm{v}^\star=0\), so
\begin{equation}
  \Gmom^\star
  =
  \Grad\left(p^\star - \Pi_{\mathrm{TG}}\right).
  \label{eq:eulerian-force}
\end{equation}
If \(p^\star=p_0+\Pi_{\mathrm{TG}}\), then
\(\Gmom^\star=0\); that is the usual unsourced incompressible
Taylor--Green balance.  If the intent is instead the pressure-canceling
Lagrangian balance stated in the MAPP description, set
\begin{equation}
  \Gmom^\star = \Grad p^\star.
  \label{eq:pressure-canceling-force}
\end{equation}
Then Eq.~\eqref{eq:primitive-momentum} gives
\begin{equation}
  \rho\frac{\dd\bm{v}}{\dd t}=0,
  \qquad
  \frac{\dd}{\dd t}\equiv\partial_t+\bm{v}\cdot\Grad,
  \label{eq:lagrangian-velocity-constant}
\end{equation}
so the velocity carried by each Lagrangian element is unchanged.

Both cases can be represented by writing the required force as a gradient,
\begin{equation}
  \Gmom^\star = \Grad\Phi.
  \label{eq:force-potential}
\end{equation}
For a pressure-canceling Lagrangian test, \(\Phi=p^\star\).  For an Eulerian
steady Taylor--Green field, \(\Phi=p^\star-\Pi_{\mathrm{TG}}\).

\section{Radiation diffusion formulation}

For a grey radiation diffusion code, the radiation flux is not an independent
evolved variable.  It is closed by Fick's law,
\begin{equation}
  \Fr = -\Ddiff \Grad \Er,
  \qquad
  \Ddiff = \frac{c}{3\rho\kappa_R},
  \label{eq:diffusion-flux}
\end{equation}
where \(\kappa_R\) is the Rosseland, or flux-mean, opacity.  The corresponding
radiation force on the gas is
\begin{equation}
  \Gmom = \frac{\rho\kappa_R}{c}\Fr
  = -\frac{1}{3}\Grad\Er.
  \label{eq:diffusion-force}
\end{equation}
The gas momentum equation is therefore equivalently
\begin{equation}
  \rho\frac{\dd\bm{v}}{\dd t} + \Grad\left(p+\frac{\Er}{3}\right)=0.
  \label{eq:diffusion-momentum}
\end{equation}

To obtain the manufactured force \(\Gmom^\star=\Grad\Phi\), choose
\begin{equation}
  \boxed{\Er^\star = \Erzero - 3\Phi.}
  \label{eq:diffusion-erad-choice}
\end{equation}
Then
\begin{equation}
  \Gmom^\star
  =
  -\frac{1}{3}\Grad\Er^\star
  =
  \Grad\Phi.
  \label{eq:diffusion-force-choice}
\end{equation}
This is the main simplification relative to a P1 two-moment code: there is no
radiation flux source to manufacture.  The diffusion closure automatically
sets
\begin{equation}
  \Fr^\star
  =
  -\Ddiff\Grad\Er^\star
  =
  3\Ddiff\Grad\Phi
  =
  \frac{c}{\rho_0\kappa_R}\Grad\Phi
  \label{eq:diffusion-flux-choice}
\end{equation}
for constant \(\rho_0\) and \(\kappa_R\).

In a Lagrangian radiation diffusion update, the comoving radiation energy
equation is commonly written
\begin{equation}
  \frac{\dd\Er}{\dd t}
  + \frac{4}{3}\Er\,\Div\bm{v}
  =
  \Div(\Ddiff\Grad\Er)
  - c\Gzero
  + \ShatE.
  \label{eq:diffusion-energy}
\end{equation}
The manufactured scalar radiation source is therefore
\begin{equation}
  \ShatE^\star
  =
  \frac{\dd\Er^\star}{\dd t}
  + \frac{4}{3}\Er^\star\Div\bm{v}^\star
  - \Div(\Ddiff\Grad\Er^\star)
  + cG^{0,\star}.
  \label{eq:diffusion-source-general}
\end{equation}
For the MAPP-style pure-Lagrange Taylor--Green setup, the desired radiation
profile is steady in the material frame, the velocity is solenoidal, and we
can choose \(G^{0,\star}=0\).  Hence
\begin{equation}
  \boxed{
  \ShatE^\star
  =
  -\Div(\Ddiff\Grad\Er^\star)
  =
  3\Div(\Ddiff\Grad\Phi).
  }
  \label{eq:diffusion-source-potential}
\end{equation}
If \(\Ddiff\) is constant, this reduces to
\begin{equation}
  \boxed{
  \ShatE^\star = 3\Ddiff\nabla^2\Phi
  = \frac{c}{\rho_0\kappa_R}\nabla^2\Phi.
  }
  \label{eq:diffusion-source-constant-d}
\end{equation}
Thus a diffusion implementation needs only this scalar radiation-energy
source.  There is no independent vector source \(\ShatF\), because the flux is
algebraically determined by Eq.~\eqref{eq:diffusion-flux}.

\section{P1 radiation fields and sources for comparison}

Choose a positive constant \(\Erzero\) and define
\begin{equation}
  \Er^\star = \Erzero - 3\Phi.
  \label{eq:erad-choice}
\end{equation}
The constant \(\Erzero\) must be large enough that \(\Er^\star>0\) everywhere.
Then the P1 radiation pressure term is
\begin{equation}
  \frac{c^2}{3}\Grad\Er^\star
  =
  -c^2\Grad\Phi
  =
  -c^2\Gmom^\star.
  \label{eq:p1-cancellation}
\end{equation}
Substitution into Eq.~\eqref{eq:general-flux-source} gives
\begin{equation}
  \ShatF^\star = \partial_t\Fr^\star.
  \label{eq:sflux-reduced}
\end{equation}
For a steady manufactured radiation flux, no external radiation-flux source is
needed:
\begin{equation}
  \boxed{\ShatF^\star=0.}
  \label{eq:sflux-zero}
\end{equation}

\subsection{Direct four-force construction}

If the problem implementation directly imposes the manufactured four-force
\((G^{0,\star},\Gmom^\star)\), one may take
\begin{equation}
  G^{0,\star}=0,\qquad
  \Fr^\star=\bm{0}.
  \label{eq:direct-force-rad-fields}
\end{equation}
Equations~\eqref{eq:general-energy-source} and \eqref{eq:sflux-zero} then give
\begin{equation}
  \boxed{\ShatE^\star=0,\qquad \ShatF^\star=\bm{0}.}
  \label{eq:direct-force-sources}
\end{equation}
In this construction the ``source'' that balances the gas is the imposed
radiation four-force itself.  The radiation moment equations remain exactly
balanced because the choice \(\Er^\star=\Erzero-3\Phi\) cancels the
\(-c^2\Gmom^\star\) term in the radiation flux equation.

\subsection{Opacity-coupled construction using a radiation energy source}

Quokka's existing single-group source update can generate a gas radiation force
through flux-mean opacity.  In the leading-order, static-frame limit
\((\beta_{\mathrm{order}}=0)\),
\begin{equation}
  \Gmom = \frac{\rho\kappa_F}{c}\Fr.
  \label{eq:opacity-force}
\end{equation}
Set \(\kappa_P=\kappa_E=0\), choose a constant
\(\kappa_F>0\), and prescribe
\begin{equation}
  \Fr^\star
  =
  \frac{c}{\rho_0\kappa_F}\Grad\Phi.
  \label{eq:flux-choice}
\end{equation}
Then Eq.~\eqref{eq:opacity-force} gives
\(\Gmom^\star=\Grad\Phi\).  With \(G^{0,\star}=0\) and steady fields,
Eq.~\eqref{eq:general-energy-source} gives the only nonzero manufactured
radiation source:
\begin{equation}
  \boxed{
  \ShatE^\star
  =
  \Div\Fr^\star
  =
  \frac{c}{\rho_0\kappa_F}\nabla^2\Phi,
  \qquad
  \ShatF^\star=\bm{0}.
  }
  \label{eq:energy-source-potential}
\end{equation}
This is the form most closely aligned with Quokka's current
\texttt{radEnergySource} interface: the pressure-balancing force is produced
by \(\Fr^\star\) and \(\kappa_F\), while the explicit radiation energy source
keeps the prescribed \(\Er^\star\) field stationary.

\section{P1 relaxation family with a diffusion limit}

Quokka has a problem-level hook for \(\ShatE\) through
\texttt{SetRadEnergySource}, but the standard update does not provide an
analogous manufactured \(\ShatF\) hook.  For a real Quokka test problem it is
therefore useful to construct a finite-parameter P1 solution using the natural
flux relaxation of the P1 equations, with \(\ShatF=0\).

For the static-frame source terms \((\beta_{\mathrm{order}}=0)\), constant
\(\rho_0\), and constant \(\kappa_R\), Quokka's reduced-speed P1 equations
are
\begin{align}
  \partial_t\Er + \Div\left(\frac{\chat}{c}\Fr\right)
  &= \SErsla, \label{eq:q-rsla-energy}\\
  \partial_t\Fr + \frac{\chat c}{3}\Grad\Er
  &= -\chat\rho_0\kappa_R\Fr. \label{eq:q-rsla-flux}
\end{align}
The gas radiation force is still
\begin{equation}
  \Gmom = \frac{\rho_0\kappa_R}{c}\Fr.
  \label{eq:q-rsla-force}
\end{equation}
Let
\begin{equation}
  \Er^\star = \Erzero - 3\Phi(\bm{x},t),
  \qquad
  \FrD = \frac{c}{\rho_0\kappa_R}\Grad\Phi.
  \label{eq:p1-diffusion-target}
\end{equation}
Then Eq.~\eqref{eq:q-rsla-flux} can be written as the relaxation equation
\begin{equation}
  \tau_F \partial_t\Fr + \Fr = \FrD,
  \qquad
  \tau_F \equiv \frac{1}{\chat\rho_0\kappa_R}.
  \label{eq:p1-relaxation}
\end{equation}
Thus the P1 asymptotic parameter is the flux relaxation time divided by the
time scale on which the diffusion flux changes.  If
\(\tau_F \partial_t \FrD\) is small, the P1 flux is
\begin{equation}
  \Fr = \FrD - \tau_F\partial_t\FrD
  + O(\tau_F^2).
  \label{eq:p1-relaxation-expansion}
\end{equation}
The gas force is then
\begin{equation}
  \Gmom
  =
  \Grad\Phi
  - \frac{1}{c\chat}\partial_t\FrD
  + O(\tau_F^2\rho_0\kappa_R/c),
  \label{eq:p1-force-error}
\end{equation}
so the pressure-canceling diffusion force is recovered as
\(\tau_F/T\to0\).  In the full-speed case this force error is
\(O(c^{-2}\partial_t\FrD)\).

\subsection{Harmonic finite-parameter solution}

For a clean analytic test, choose a harmonic potential
\begin{equation}
  \Phi(\bm{x},t) = \operatorname{Re}\left[\widehat{\Phi}(\bm{x})e^{i\omega t}\right].
  \label{eq:harmonic-potential}
\end{equation}
After the homogeneous relaxation transient has decayed, the exact P1 flux with
\(\ShatF=0\) is
\begin{equation}
  \Fr^\tau
  =
  \frac{c}{\rho_0\kappa_R}
  \Grad
  \operatorname{Re}\left[
    \frac{\widehat{\Phi}(\bm{x})e^{i\omega t}}{1+i\omega\tau_F}
  \right].
  \label{eq:harmonic-flux}
\end{equation}
Consequently
\begin{equation}
  \Gmom^\tau
  =
  \Grad\Psi^\tau,
  \qquad
  \Psi^\tau
  \equiv
  \operatorname{Re}\left[
    \frac{\widehat{\Phi}(\bm{x})e^{i\omega t}}{1+i\omega\tau_F}
  \right].
  \label{eq:harmonic-force-potential}
\end{equation}
The finite-\(\tau_F\) gas pressure that is exactly balanced by the P1
radiation force is therefore \(p^\tau=p_0+\Psi^\tau\) for the Lagrangian
pressure-canceling version, or
\(p^\tau=p_0+\Pi_{\mathrm{TG}}+\Psi^\tau\) for the Eulerian steady
Taylor--Green velocity field.  The constant \(p_0\) does not affect the force;
it is included to keep the gas pressure positive.  The diffusion pressure is
recovered when \(\tau_F\to0\).

The relative amplitude error in the force potential is
\begin{equation}
  \frac{\left|\Psi^\tau-\Phi\right|}{|\Phi|}
  =
  \frac{\omega\tau_F}{\sqrt{1+(\omega\tau_F)^2}}
  =
  \omega\tau_F + O((\omega\tau_F)^2).
  \label{eq:potential-error}
\end{equation}
The same relative error applies to the radiation flux and radiation force for
each Fourier mode.  This gives a direct analytic estimate for the modeling
error caused by a finite asymptotic parameter.

The exact stored radiation-energy source for Quokka's reduced-speed equations
is
\begin{equation}
  \SErsla^\tau
  =
  \partial_t\Er^\star
  + \Div\left(\frac{\chat}{c}\Fr^\tau\right)
  =
  -3\partial_t\Phi
  + \frac{\chat}{\rho_0\kappa_R}\nabla^2\Psi^\tau.
  \label{eq:finite-tau-stored-source}
\end{equation}
For thermal radiation groups, Quokka internally multiplies the user-provided
\texttt{radEnergySource} by \(\chat/c\).  Therefore the source passed through
\texttt{SetRadEnergySource} should be
\begin{equation}
  S_{\mathrm{input}}^\tau
  =
  \frac{c}{\chat}\SErsla^\tau
  =
  -3\frac{c}{\chat}\partial_t\Phi
  + \frac{c}{\rho_0\kappa_R}\nabla^2\Psi^\tau.
  \label{eq:finite-tau-input-source}
\end{equation}
In the diffusion limit, \(\Psi^\tau\to\Phi\), and this reduces to the
appropriate reduced-speed diffusion source for the same target
\(\Er^\star=\Erzero-3\Phi\).

\subsection{Implications for a Quokka test}

A stationary manufactured profile has \(\omega=0\), and then
\(\Psi^\tau=\Phi\): the P1 construction is exactly balanced for any positive
\(\kappa_R\).  Such a test is useful for checking the source implementation,
but it cannot measure the asymptotic P1-to-diffusion error.

To test the finite asymptotic parameter, use a time-dependent potential such
as
\begin{equation}
  \Phi(\bm{x},t)
  =
  A_p \cos(\omega t)
  \left[\cos(2kx)+\cos(2ky)\right].
  \label{eq:test-harmonic-potential}
\end{equation}
For this mode, Eq.~\eqref{eq:potential-error} predicts the finite-\(\tau_F\)
pressure, flux, and radiation-force error.  Since
\(\nabla^2[\cos(2kx)+\cos(2ky)]=-4k^2[\cos(2kx)+\cos(2ky)]\), the amplitude of
the finite-\(\tau_F\) correction to the diffusion energy source is
\begin{equation}
  12D_{\chat}k^2A_p
  \frac{\omega\tau_F}{\sqrt{1+(\omega\tau_F)^2}},
  \qquad
  D_{\chat}\equiv\frac{\chat}{3\rho_0\kappa_R}.
  \label{eq:source-error-amplitude}
\end{equation}
This provides a closed-form tolerance target: after spatial and temporal
discretization errors are made small, the remaining difference between a
finite-\(\tau_F\) P1 calculation and the diffusion manufactured solution should
scale linearly with \(\omega\tau_F\).

\subsection{Practical parameter choices}

A practical first implementation should keep the gas nearly incompressible,
the radiation flux safely subluminal, and the flux relaxation error large
enough to measure above discretization error.  The following dimensionless
choice is a useful starting point:
\begin{center}
\begin{tabular}{ll}
\hline
quantity & recommended value \\
\hline
domain & \([0,1]^2\), periodic \\
Taylor--Green wavenumber & \(k=2\pi\) \\
pressure/source mode & \(q=\cos(4\pi x)+\cos(4\pi y)\) \\
density & \(\rho_0=1\) \\
gas adiabatic index & \(\gamma=5/3\) \\
gas pressure floor/base & \(p_0=0.1\) \\
pressure perturbation & \(\Phi=A_p q\cos(\omega t)\), \(A_p=5\times10^{-3}\) \\
radiation energy offset & \(\Erzero=1\) \\
physical light speed & \(c=100\) \\
reduced light speed & \(\chat=1\) \\
velocity amplitude & \(U=10^{-2}\) or smaller \\
radiation source terms & \(\kappa_P=\kappa_E=0\), \(\kappa_F=\kappa_R\) \\
four-force order & \(\beta_{\mathrm{order}}=0\) \\
driving frequency & \(\omega=5\) \\
default opacity & \(\kappa_R=100\) \\
resolution & \(64^2\) for development, \(128^2\) for convergence \\
comparison time & \(t_f=\pi/(2\omega)\) \\
\hline
\end{tabular}
\end{center}
With these values,
\begin{equation}
  \Omega \equiv \omega\tau_F
  = \frac{\omega}{\chat\rho_0\kappa_R}
  = 0.05,
  \label{eq:recommended-omega}
\end{equation}
so the finite-P1 force and flux differ from the diffusion manufactured
solution by approximately \(5\%\).  This is large enough to measure, but still
small enough that the solution is visibly in the diffusion regime.  At
\(64^2\), the cell optical depth is
\begin{equation}
  \tau_{\mathrm{cell}}=\rho_0\kappa_R\Delta x = \frac{100}{64}\approx1.56,
  \label{eq:recommended-cell-tau}
\end{equation}
and the optical depth across one pressure-mode wavelength is
\begin{equation}
  \tau_\lambda=\rho_0\kappa_R\left(\frac{1}{2}\right)=50.
  \label{eq:recommended-wavelength-tau}
\end{equation}
The maximum reduced flux is approximately
\begin{equation}
  f_{\max}
  =
  \frac{|\Fr|}{c\Er}
  \sim
  \frac{4\pi A_p\sqrt{2}}{\rho_0\kappa_R(\Erzero-6A_p)}
  \approx 9\times10^{-4},
  \label{eq:recommended-reduced-flux}
\end{equation}
well inside the P1/M1 admissible region.

For an asymptotic sweep, keep \(\omega=5\) and use
\begin{equation}
  \kappa_R \in \{50,100,200,400\},
  \qquad
  \Omega \in \{0.1,0.05,0.025,0.0125\}.
  \label{eq:recommended-sweep}
\end{equation}
The finite-asymptotic error should decrease linearly with \(\Omega\) once the
mesh and timestep errors are smaller than this modeling error.  The quarter
period comparison time \(t_f=\pi/(2\omega)\) is preferable to a full period
because it exposes the \(O(\Omega)\) phase lag directly; at a full period the
leading phase error cancels in a final-time snapshot.

These parameters assume the problem initializes the exact finite-\(\tau_F\)
flux in Eq.~\eqref{eq:harmonic-flux}, rather than the diffusion flux
\(\FrD\).  That removes the homogeneous flux-relaxation transient.  If the
gas pressure is made explicitly time dependent to match
\(\Psi^\tau\), an ideal-gas implementation also needs a matching gas internal
energy source,
\begin{equation}
  S_{\mathrm{gas},e}
  =
  \frac{1}{\gamma-1}
  \left(\partial_t p^\tau+\bm{v}\cdot\Grad p^\tau\right),
  \label{eq:recommended-gas-energy-source}
\end{equation}
or else the test should first be run with \(U=0\) to isolate the radiation
asymptotics.

\subsubsection{Gas and radiation positivity}

Pressure positivity should be enforced analytically by choosing \(p_0\) larger
than the largest possible negative pressure perturbation.  In the Lagrangian
pressure-canceling version,
\begin{equation}
  p^\tau = p_0 + \Psi^\tau,
  \qquad
  |\Psi^\tau|
  \le
  \frac{2A_p}{\sqrt{1+(\omega\tau_F)^2}},
  \label{eq:lagrangian-pressure-bound}
\end{equation}
because \(|\cos(4\pi x)+\cos(4\pi y)|\le2\).  Hence a sufficient condition is
\begin{equation}
  \boxed{
  p_0 >
  \frac{2A_p}{\sqrt{1+(\omega\tau_F)^2}}.
  }
  \label{eq:lagrangian-pressure-positivity}
\end{equation}
For the Eulerian steady Taylor--Green velocity field, the gas pressure also
contains \(\Pi_{\mathrm{TG}}\).  With \(k=2\pi\),
\begin{equation}
  \Pi_{\mathrm{TG}}
  =
  \frac{\rho_0U^2}{4}
  \left[\cos(4\pi x)+\cos(4\pi y)\right],
  \label{eq:practical-pitg}
\end{equation}
so a sufficient pressure-positivity condition is
\begin{equation}
  \boxed{
  p_0 >
  2\left[
    \frac{\rho_0U^2}{4}
    +
    \frac{A_p}{\sqrt{1+(\omega\tau_F)^2}}
  \right].
  }
  \label{eq:eulerian-pressure-positivity}
\end{equation}
For the recommended values \(p_0=0.1\), \(A_p=5\times10^{-3}\),
\(\rho_0=1\), \(U=10^{-2}\), and \(\omega\tau_F=0.05\), this gives
\begin{equation}
  p_{\min}
  \gtrsim
  0.1 - 2\left(2.5\times10^{-5}+4.99\times10^{-3}\right)
  \approx 8.995\times10^{-2},
  \label{eq:recommended-pressure-min}
\end{equation}
so the gas pressure remains comfortably positive.

Radiation energy positivity is separate.  With the practical convention that
\(\Phi\) is the non-constant pressure potential and \(p_0\) is a gas-pressure
offset,
\begin{equation}
  \Er^\star=\Erzero-3\Phi,
  \qquad
  E_{\mathrm r,\min}^\star
  \ge
  \Erzero-6A_p.
  \label{eq:radiation-energy-bound}
\end{equation}
Thus \(\Erzero>6A_p\) is sufficient for the Lagrangian pressure-canceling
version.  The recommended \(\Erzero=1\) and \(A_p=5\times10^{-3}\) give
\(E_{\mathrm r,\min}^\star\ge0.97\).

\section{Explicit pressure-canceling Taylor--Green source}

Let the pressure to be canceled be
\begin{equation}
  p^\star(x,y)
  =
  p_0 + A_p\left[\cos(2kx)+\cos(2ky)\right].
  \label{eq:generic-pressure}
\end{equation}
For the MAPP-style pressure-canceling Lagrangian interpretation,
\(\Phi=p^\star\).  Then
\begin{align}
  \Grad p^\star
  &=
  -2kA_p
  \begin{pmatrix}
    \sin(2kx)\\
    \sin(2ky)\\
    0
  \end{pmatrix},
  \label{eq:pressure-gradient}\\
  \nabla^2 p^\star
  &=
  -4k^2 A_p\left[\cos(2kx)+\cos(2ky)\right].
  \label{eq:pressure-laplacian}
\end{align}
The diffusion manufactured radiation field, algebraic flux, and scalar source
are
\begin{align}
  \Er^\star
  &=
  \Erzero - 3p^\star,
  \label{eq:pressure-erad}\\
  \Fr^\star
  &=
  -\frac{2ckA_p}{\rho_0\kappa_R}
  \begin{pmatrix}
    \sin(2kx)\\
    \sin(2ky)\\
    0
  \end{pmatrix},
  \label{eq:pressure-frad}\\
  \ShatE^\star
  &=
  -\frac{4ck^2A_p}{\rho_0\kappa_R}
  \left[\cos(2kx)+\cos(2ky)\right].
  \label{eq:pressure-source}
\end{align}
There is no separate \(\ShatF^\star\) in a diffusion formulation because
\(\Fr^\star\) is defined algebraically by Eq.~\eqref{eq:diffusion-flux}.
If \(A_p=\rho_0U^2/4\), this becomes
\begin{align}
  \Fr^\star
  &=
  -\frac{cU^2k}{2\kappa_R}
  \begin{pmatrix}
    \sin(2kx)\\
    \sin(2ky)\\
    0
  \end{pmatrix},
  \label{eq:canonical-frad}\\
  \ShatE^\star
  &=
  -\frac{cU^2k^2}{\kappa_R}
  \left[\cos(2kx)+\cos(2ky)\right].
  \label{eq:canonical-source}
\end{align}
The leading-order opacity-coupled P1 construction gives the same scalar
source after replacing \(\kappa_R\) by \(\kappa_F\), but P1 has an independent
radiation flux variable and therefore requires the separate balance discussed
above.

For a strictly Eulerian steady Taylor--Green field, use instead
\(\Phi=p^\star-\Pi_{\mathrm{TG}}\).  With the pressure in
Eq.~\eqref{eq:generic-pressure},
\begin{equation}
  \Phi
  =
  p_0 +
  \left(A_p-\frac{\rho_0U^2}{4}\right)
  \left[\cos(2kx)+\cos(2ky)\right],
  \label{eq:eulerian-potential-explicit}
\end{equation}
and Eq.~\eqref{eq:diffusion-source-constant-d} applies with \(A_p\) replaced by
\(A_p-\rho_0U^2/4\).  In the canonical Taylor--Green case
\(A_p=\rho_0U^2/4\), \(\Phi\) is constant and the radiation force and
manufactured radiation energy source both vanish.

\section{Why the density remains constant}

The mass equation can be written in material form as
\begin{equation}
  \frac{\dd\rho}{\dd t} = -\rho\,\Div\bm{v}.
  \label{eq:material-mass}
\end{equation}
For the prescribed Taylor--Green velocity,
\(\Div\bm{v}_{\mathrm{TG}}=0\) by Eq.~\eqref{eq:tgv-divergence}.  Therefore
\begin{equation}
  \frac{\dd\rho}{\dd t}=0.
  \label{eq:density-material-constant}
\end{equation}
If \(\rho=\rho_0\) initially, then \(\rho=\rho_0\) for all time.  Equivalently,
in Eulerian form,
\begin{equation}
  \partial_t\rho^\star + \Div(\rho^\star\bm{v}_{\mathrm{TG}})
  =
  0+\rho_0\Div\bm{v}_{\mathrm{TG}}
  =
  0.
  \label{eq:density-eulerian-constant}
\end{equation}
This proves the constant-density part of the manufactured solution.

\section{Why the velocity remains unchanged}

In the diffusion formulation, the pressure-canceling construction is immediate.
With \(\Phi=p^\star\),
\begin{equation}
  p^\star + \frac{\Er^\star}{3}
  =
  p^\star + \frac{\Erzero-3p^\star}{3}
  =
  \frac{\Erzero}{3},
  \label{eq:total-pressure-constant}
\end{equation}
which is spatially constant.  The diffusion momentum equation
\eqref{eq:diffusion-momentum} therefore gives
\begin{equation}
  \rho\frac{\dd\bm{v}}{\dd t}
  =
  -\Grad\left(p^\star+\frac{\Er^\star}{3}\right)
  =
  \bm{0}.
  \label{eq:velocity-diffusion-proof}
\end{equation}
Thus the velocity attached to each Lagrangian fluid element remains unchanged,
as specified in the MAPP description.

For the pressure-canceling Lagrangian construction, the manufactured radiation
force is \(\Gmom^\star=\Grad p^\star\).  Substitution into the primitive
momentum equation gives
\begin{equation}
  \rho\frac{\dd\bm{v}}{\dd t}+\Grad p^\star=\Grad p^\star,
  \qquad
  \frac{\dd\bm{v}}{\dd t}=\bm{0}.
  \label{eq:velocity-lagrangian-proof}
\end{equation}
This is the same statement written in Quokka's four-force sign convention.

For an Eulerian code test in which the exact velocity field is required to be
the same function of position at all times,
\(\bm{v}^\star(x,y,t)=\bm{v}_{\mathrm{TG}}(x,y)\), use the more general force
in Eq.~\eqref{eq:eulerian-force}.  Then
\begin{align}
  \rho_0\partial_t\bm{v}^\star
  &=
  \Gmom^\star
  - \rho_0(\bm{v}^\star\cdot\Grad)\bm{v}^\star
  - \Grad p^\star \nonumber\\
  &=
  \Grad(p^\star-\Pi_{\mathrm{TG}})
  - \left[-\Grad\Pi_{\mathrm{TG}}\right]
  - \Grad p^\star \nonumber\\
  &= \bm{0}.
  \label{eq:velocity-eulerian-proof}
\end{align}
This proves that the Eulerian velocity field is stationary when the
manufactured force includes both the hydrodynamic pressure residual and the
Taylor--Green advective acceleration.

\section{Reduced-speed-of-light form used by Quokka}

When Quokka uses a reduced speed of light \(\chat\ne c\), the radiation
transport terms in the stored radiation variables are scaled as
\begin{align}
  \partial_t\Er + \Div\left(\frac{\chat}{c}\Fr\right)
  &= -\chat\Gzero + \SErsla,
  \label{eq:rsla-energy}\\
  \partial_t\Fr + \Div(\chat c\Prad)
  &= -\chat c\Gmom + \SFrsla.
  \label{eq:rsla-flux}
\end{align}
Under P1,
\begin{align}
  \SErsla^\star
  &=
  \partial_t\Er^\star
  + \Div\left(\frac{\chat}{c}\Fr^\star\right)
  + \chat G^{0,\star},
  \label{eq:rsla-energy-source}\\
  \SFrsla^\star
  &=
  \partial_t\Fr^\star
  + \frac{\chat c}{3}\Grad\Er^\star
  + \chat c\Gmom^\star.
  \label{eq:rsla-flux-source}
\end{align}
The same choice \(\Er^\star=\Erzero-3\Phi\) makes
\(\SFrsla^\star=0\) for steady fields.  For the opacity-coupled
construction, \(\Fr^\star=c\Grad\Phi/(\rho_0\kappa_F)\) still gives the gas
force \(\Gmom^\star=\rho_0\kappa_F\Fr^\star/c=\Grad\Phi\), and the stored
energy source becomes
\begin{equation}
  \SErsla^\star
  =
  \frac{\chat}{\rho_0\kappa_F}\nabla^2\Phi.
  \label{eq:rsla-energy-source-final}
\end{equation}
Quokka's thermal-group \texttt{radEnergySource} input is a physical luminosity
density and is internally multiplied by \(\chat/c\), so the value passed
through that interface should be the full-speed expression
\(\Div\Fr^\star=c\nabla^2\Phi/(\rho_0\kappa_F)\).

\end{document}
